\section{A Finite-Dimensional Real Gleason-Type Theorem}

\begin{theorem}[Real effect-space theorem]
Fix $d\ge 2$. Let $\Eff_d$ be the set of all real symmetric $d\times d$ matrices $E$ satisfying $0\le E\le I$. Suppose
\[
f:\Eff_d\to [0,1]
\]
satisfies
\[
f(I)=1
\]
and
\[
f(E+F)=f(E)+f(F)
\]
whenever $E,F,E+F\in \Eff_d$. Then there exists a unique real symmetric positive semidefinite matrix $\rho$ with $\mathrm{tr}(\rho)=1$ such that
\[
f(E)=\mathrm{tr}(\rho E)
\qquad\text{for every }E\in \Eff_d.
\]
\end{theorem}

\begin{proof}
We break the proof into four simple steps.

\paragraph{Step 1: homogeneity on each ray.}
Fix an effect $E$. The map
\[
h_E(t)=f(tE),\qquad t\in[0,1],
\]
is bounded, and it satisfies
\[
h_E(x+y)=h_E(x)+h_E(y)
\]
whenever $x+y\le 1$, because then $(x+y)E$, $xE$, and $yE$ all remain effects. By the interval lemma,
\[
f(tE)=t\,f(E)\qquad (0\le t\le 1).
\]

\paragraph{Step 2: extend from effects to all positive semidefinite matrices.}
Let $A$ be any positive semidefinite real symmetric matrix. Choose a positive number $s$ large enough that $A/s\in\Eff_d$; for instance, any $s$ at least as large as the largest eigenvalue of $A$ will do. Define
\[
F(A)=s\,f(A/s).
\]
This does not depend on the chosen $s$. Indeed, if $A/s$ and $A/t$ are both effects and $s\le t$, then
\[
A/t=(s/t)(A/s),
\]
so Step 1 gives
\[
t\,f(A/t)=t\,(s/t)\,f(A/s)=s\,f(A/s).
\]
Thus $F$ is well defined.

Now let $A,B$ be positive semidefinite. Choose $u$ large enough that $(A+B)/u$ is an effect. Then
\[
F(A+B)=u\,f((A+B)/u)=u\,f(A/u)+u\,f(B/u)=F(A)+F(B).
\]
So $F$ is additive on the positive semidefinite cone.

\paragraph{Step 3: extend to all symmetric matrices.}
Every real symmetric matrix $H$ can be written as
\[
H=A-B
\]
with $A,B$ positive semidefinite. Define
\[
L(H)=F(A)-F(B).
\]
This is well defined. If $H=C-D$ also, then $A+D=B+C$, and additivity of $F$ gives
\[
F(A)+F(D)=F(B)+F(C),
\]
so $F(A)-F(B)=F(C)-F(D)$.

Linearity of $L$ is now immediate:
\[
L((A-B)+(C-D))=F(A+C)-F(B+D)=L(A-B)+L(C-D),
\]
and $L(\lambda H)=\lambda L(H)$ follows from the positive-scalar homogeneity of $F$.

\paragraph{Step 4: write $L$ as a trace rule.}
Let $e_1,\dots,e_d$ be the standard basis vectors. Define
\[
P_i=e_i e_i^\top
\]
and, for $i<j$,
\[
S_{ij}=\tfrac12(e_i+e_j)(e_i+e_j)^\top.
\]
Each $P_i$ and $S_{ij}$ is an effect; in fact, $S_{ij}$ is the rank-one projection onto the unit vector $(e_i+e_j)/\sqrt2$. Set
\[
\rho_{ii}=f(P_i),\qquad
\rho_{ij}=f(S_{ij})-\tfrac12 f(P_i)-\tfrac12 f(P_j)\quad (i<j),
\]
and define $\rho_{ji}=\rho_{ij}$. This gives a real symmetric matrix $\rho$.

Now let
\[
X_{ij}=\tfrac12(E_{ij}+E_{ji})=S_{ij}-\tfrac12 P_i-\tfrac12 P_j.
\]
So $X_{ij}$ is exactly the symmetric off-diagonal direction recovered from values of $f$ on honest effects.
By construction,
\[
L(P_i)=\rho_{ii},\qquad L(X_{ij})=\rho_{ij}.
\]
Every real symmetric matrix $H=(h_{ij})$ has the basis expansion
\[
H=\sum_{i=1}^d h_{ii}P_i+2\sum_{1\le i<j\le d} h_{ij}X_{ij}.
\]
Therefore,
\[
L(H)=\sum_{i=1}^d h_{ii}\rho_{ii}+2\sum_{1\le i<j\le d} h_{ij}\rho_{ij}
=\mathrm{tr}(\rho H).
\]

We still need positivity and normalization. If $v$ is any vector, then $vv^\top$ is positive semidefinite, so
\[
v^\top \rho v=\mathrm{tr}(\rho vv^\top)=L(vv^\top)=F(vv^\top)\ge 0.
\]
Hence $\rho$ is positive semidefinite.

Also,
\[
I=P_1+\cdots+P_d,
\]
and the partial sums remain effects, so
\[
\mathrm{tr}(\rho)=\sum_{i=1}^d \rho_{ii}
=\sum_{i=1}^d f(P_i)=f(I)=1.
\]
Finally, if $E\in\Eff_d$, then $E$ itself is positive semidefinite, so
\[
f(E)=F(E)=L(E)=\mathrm{tr}(\rho E).
\]
Uniqueness is immediate because the values on the basis effects determine every entry of $\rho$.
\end{proof}

The argument above is intentionally concrete. No Hilbert-space lattice, no measure theory, and no quantum postulates are needed. The only real inputs are bounded additivity on $[0,1]$, positive semidefinite decompositions, and a basis expansion of symmetric matrices.
